\section{Interpretação}
As branches mostram três perguntas diferentes: se o mecanismo pode executar corretamente num corte live, se altera o comportamento numérico em integração vertical, e se reduz o caminho crítico de um workload. A resposta positiva observada para logits após reinjection/split-graph só cobre aqueles estados capturados; a divergência multi-layer e o contrato N2 continuam válidos no seu próprio âmbito \historicalclaim{CL-S02-REINJECTION}{} \historicalclaim{CL-S02-SPLIT}{} \historicalclaim{CL-S02-VERTICAL}{} \historicalclaim{CL-S02-N2}{}.

A curva CPU-MoE e o resultado async são diagnósticos úteis de trade-off/crossover, mas não demonstram um winner repetível, front de Pareto estável, benefício de modelo inteiro ou tráfego físico. O gate prospectivo que segue a curva é uma medição de CPU-vs-transfer-vs-GPU; o resultado async requer confirmação de cauda antes de uma integração ampla \historicalclaim{CL-S02-NCPU-PILOT}{} \historicalclaim{CL-S02-NCPU-CONTINUUM}{} \historicalclaim{CL-S02-ASYNC}{}.

Este limite é compatível com a literatura: colocação CPU/GPU, caching de experts e overlap/pipeline já têm sistemas publicados. As referências e comparadores mais próximos — incluindo MoE-Infinity, PowerInfer, MoE-Lightning, Klotski e DALI — estão ligados às claims de mecanismo em R01 \cite{xue2024moeinfinity,song2024powerinfer,cao2025moelightning,fang2025klotski,zhu2026dali}. S02 não reclama novidade para nenhum desses mecanismos.
